\documentclass{beamer}
\usepackage[utf8]{inputenc}
\usepackage[italian]{babel}
\usepackage[backend=biber,style=alphabetic,hyperref,backref]{biblatex}
\usepackage{listings}
% ---- IMPOSTAZIONI PACCHETTI ----
\definecolor{forestgreen}{rgb}{0.13, 0.55, 0.13}

\lstset{
    language=[LaTeX]Tex,%C++,
    keywordstyle=\color{blue}, %\bfseries,
    basicstyle=\small\ttfamily,
    commentstyle=\color{forestgreen}\ttfamily,
    stringstyle=\rmfamily,
    numbers=left,%
    numberstyle=\scriptsize, %\tiny
    stepnumber=1,
    numbersep=8pt,
    showstringspaces=false,
    breaklines=true,
    frameround=ftff,
    frame=single
}
\usepackage{../UnipdTheme/Padova/beamerthemePadova}
\usepackage[absolute,overlay]{textpos}
\usepackage[export]{adjustbox}

\newcommand{\myWiki}[2][] {
	\href{http://en.wikipedia.org/wiki/#2}
	{\ifthenelse{\equal{#1}{}}{http://en.wikipedia.org/wiki/#2}{#1}}
}

\newcommand{\myConcatOne}[2] {
	#1#2
}

\newcommand{\myConcatTwo}[2][] {
	{\ifthenelse{\equal{#1}{}}{#2}{#1#2}}
}

\bibliography{bibliografia}

\title{Bibliografia e creazione dei comandi}
\subtitle{Solo per stomaci forti}
\author{Davide Polonio, Marco Zanella \& Emanuele Carraro}
\date{AA 2017-2018}


\begin{document}

	\maketitle
	\begin{frame}{Bibliografia}

Spesso nei documenti come le tesi si vuole fare riferimento a testi esterni o
articoli vari 

\vfill

Ci sono 2 metodi principali
\begin{itemize}
	\item manuale
	\item automatica
\end{itemize}

\vfill

La scelta dipende da quante informazioni si vogliono o devono aggiungere alla
bibliografia e dalla personalizzazione che si vuole dare alle citazioni

\begin{figure}
	\centering
	\includegraphics[scale=0.25]{res/images/bibliografia}
\end{figure}

\end{frame}
	\begin{frame}{Bibliografia manuale}

\begin{textblock*}{2cm}(10cm,1.5cm)
      \includegraphics[scale=0.30]{res/images/manual}
\end{textblock*}

È la scelta più semplice ma presenta alcune limitazioni

\vfill

Comandi da utilizzare:
\begin{itemize}
	\item \texttt{\textbackslash{}begin\{thebibliography\}\{n\}\dots{}
	\textbackslash{}end\{thebibliography\}} ambiente in cui dichiarare le opere
	\item \texttt{\textbackslash{}bibitem[etichetta\_personalizzate]
	\{identificativo\}} dove
	\begin{itemize}
		\item \textcolor{red}{\texttt{etichetta\_personalizzate}} è opzionale
		e permette di definire etichette personalizzate
		\item \textcolor{red}{\texttt{identificativo}} è una label e si
		consiglia la sintassi \texttt{autore:titolo}
	\end{itemize}
	\item \texttt{\textbackslash{}cite\{identificativo\}} per inserire una
	citazione
\end{itemize}

\end{frame}
    \begin{frame}[fragile]{Inserimento bibliografia manuale}

\begin{textblock*}{2cm}(10cm,1.5cm)
      \includegraphics[scale=0.30]{res/images/manual}
\end{textblock*}

Nel caso di \texttt{book} o \texttt{report}
\begin{lstlisting}
\cleardoublepage
%\phantomsection
\addcontentsline{toc}{chapter}{\bibname}
\end{lstlisting}

\vfill

Nel caso di \texttt{article}
\begin{lstlisting}
\clearpage
%\phantomsection
\addcontentsline{toc}{section}{\refname}
\end{lstlisting}

\end{frame}
    \begin{frame}[fragile]{Esempio bibliografia manuale - 1}

\begin{exampleblock}{Bibliografia manuale - Preambolo, testo e inserimento}
	\begin{lstlisting}
\documentclass{book}

\begin{document}

\dots{}citazione di un'opera~\cite{autore1:titolo1} bla bla bla\dots{}
\\
\dots{}citazione di un'opera~\cite{autore2:titolo2} bla bla bla\dots{}

\cleardoublepage
%\phantomsection
\addcontentsline{toc}{chapter}{\bibname}
	\end{lstlisting}
\end{exampleblock}

\end{frame}
    \begin{frame}[fragile]{Esempio bibliografia manuale - 2}

\begin{exampleblock} {Bibliografia manuale - codice bibliografia}
	\begin{lstlisting}
\begin{thebibliography}{9}

\bibitem{autore1:titolo1} Cognome, Nome(1979), \emph{Titolo completo}.
\bibitem[etichetta\_personalizzata]{autore2:titolo2} Cognome, Nome(1979),
\emph {Titolo completo}.

\end{thebibliography}

\end{document}
	\end{lstlisting}
\end{exampleblock}

\end{frame}
    \begin{frame}{Pro e contro bibliografia manuale}

\begin{textblock*}{2cm}(10.7cm,2cm)
      \includegraphics[scale=0.25]{res/images/pro}
\end{textblock*}

Vantaggi
\begin{itemize}
	\item Molto semplice
	\item Adatta a documenti brevi con poche citazioni
\end{itemize}

\vfill

\begin{textblock*}{2cm}(10.7cm,5cm)
      \includegraphics[scale=0.25]{res/images/contro}
\end{textblock*}

Svantaggi
\begin{itemize}
	\item Opere non riordinate
	\item Deve essere riscritta per ogni documento
	\item Per cambiare lo stile si deve mettere mano ad ogni singola voce
\end{itemize}

\end{frame}
    \begin{frame}{Bibliografia automatica}

\begin{textblock*}{2cm}(10cm,2cm)
      \includegraphics[scale=0.28]{res/images/automatic}
\end{textblock*}

È più complessa ma permette un grado di personalizzazione maggiore

\vfill

Per utilizzarla è necessario
\begin{itemize}
	\item Utilizzare il pacchetto \texttt{biblatex} al quale è possibile è possibile specificare una serie di opzioni
	\begin{itemize}
		\item \texttt{backend=biber}
		\item \texttt{style} per definire lo stile delle citazioni
		\item \texttt{hyperref} e \texttt{backref} nel caso in cui si voglia
		che ci siano nel documento i collegamenti tra citazione e dati
		bibliografici
	\end{itemize}
	\item importare il pacchetto per le citazioni con il comando 
	\texttt{\textbackslash{}usepackage[autostyle,italian=guillemets]
	\{csquotes\}}
	\item È necessario creare un \emph{database} delle opere
\end{itemize}

\end{frame}
    \begin{frame}{Bibliografia automatica - Stili di citazione}

\begin{textblock*}{2cm}(10cm,2cm)
      \includegraphics[scale=0.28]{res/images/automatic}
\end{textblock*}

\begin{description}
	\item[\textbf{numeric}] Riferimento numerico es. [1], [2], \dots{}
	\item[\textbf{alphabetic}] Riferimento misto es. [Mor07]
	\item[\textbf{authoryear}] Autore ed anno es. [Mori,2007]
	\item[\textbf{authortitle}] Autore e titolo dell'opera es. [Mori,Opera]
\end{description}

\end{frame}
    \begin{frame}[fragile]{Bibliografia automatica - Database opere}

\begin{exampleblock}{Esempio opera}
	\begin{lstlisting}
@book{eco:tesi,
  author     = {Eco, Umberto},
  title      = {Come si fa una tesi di laurea},
  publisher  = {Bompiani},
  date       = {1977},
  location   = {Milano},
}
	\end{lstlisting}
\end{exampleblock}

\end{frame}
    \begin{frame}[fragile]{Inserimento bibliografia automatica}

\begin{lstlisting}
%vedi bibliografia manuale per il significato
\cleardoublepage
%per inserire anche i riferimenti senza citazioni
\nocite{*}
%per stampare effettivamente la bibliografia
\printbibliography
\end{lstlisting}

\end{frame}
    \begin{frame}[fragile]{Esempio bibliografia automatica - 1}

\begin{exampleblock}{Bibliografia automatica - Preambolo, testo e inserimento}
	\begin{lstlisting}
\documentclass{book}
\usepackage[utf8]{inputenc}
\usepackage{hyperref}
\usepackage[autostyle, italian=guillemets]{csquotes}
\usepackage[backend=biber, style=alphabetic, hyperref, backref]{biblatex}
\addbibresource{bibliografia.bib}
\begin{document}
\dots{}citazione di un'opera~\cite{eco:tesi} bla bla bla\dots{}\\
\dots{}citazione di un'opera~\cite{mori:tesi} bla bla bla\dots{}
\cleardoublepage \nocite {*}
\printbibliography
\end{document}
	\end{lstlisting}
\end{exampleblock}

\end{frame}
    \begin{frame}[fragile]{Esempio bibliografia automatica - 2}

\begin{exampleblock} {File bibliografia.bib}
	\begin{lstlisting}
@book{eco:tesi,
  author       = {Eco, Umberto},
  title        = {Come si fa una tesi di laurea},
  publisher    = {Bompiani},
  date         = {1977},
  location     = {Milano},
}
@article{mori:tesi,
  author       = {Mori, Lapo Filippo},
  title        = {Scrivere la tesi di laurea con \LaTeXe},
  journaltitle = {Giornale},
  number       = {3},
  date         = {2007},
}
	\end{lstlisting}
\end{exampleblock}

\end{frame}
    \begin{frame}[fragile]{Esercizio bibliografia}

\begin{block}{Esercizio bibliografia}
    
Definire la seguente lista di citazioni.
Per la prima citazione, usate \texttt{\textbackslash citet}, per la seconda \texttt{\textbackslash citep} e nella terza ricordatevi di specificare la pagina.

\begin{enumerate}
\item ``La vita è come una commedia, non importa quanto è lunga, ma come è recitata.'' \citet{seneca}
\item ``Ciò che non mi distrugge mi rende più forte.'' \citep{nietzsche}
\item ``Tutti gli esseri della Flatlandia presentano al nostro ``occhio'' il medesimo, o quasi il medesimo aspetto, quello cioè di una Linea Retta.'' \citep[p. 42]{abbot}
\end{enumerate}

\end{block}

\end{frame}

    \begin{frame}[fragile]{Soluzione bibliografia - 1/2}

\begin{block}{Lista}
    
\begin{lstlisting}
    \usepackage{natbib}
    ...
    \begin{enumerate}
      \item \dots \citet{seneca}
      \item \dots \citep{nietzsche}
      \item \dots \citep[p. 42]{abbot}
    \end{enumerate}
    \clearpage
    \bibliographystyle{plainnat}
    \bibliography{bibliografia}
\end{lstlisting}

\end{block}

\end{frame}

    \begin{frame}[fragile]{Soluzione bibliografia - 2/2}

\begin{block}{Bibliografia}

\begin{lstlisting}[basicstyle=\tiny]
@book{seneca,
    author = {Seneca},
    year = {45},
    title = {Opere di Seneca},
    publisher = {Mondadori},
}
@book{nietzsche,
    author = {Nietzsche},
    year = {1888},
    title = {Ecce homo},
    publisher = {Feltrinelli},
}
@book{abbot,
    author = {Edwin Abbott Abbott},
    year = {2003},
    title = {Flatlandia},
    publisher = {Adelphi},
}
\end{lstlisting}

\end{block}

\end{frame}

    \begin{frame}[fragile]{Creazione comandi personalizzati}

Spesso può risultare utile creare un nuovo comando
\begin{itemize}
	\item Per eseguire ``azioni'' ripetitive
	\item In modo tale da modificare in modo semplice più parti
\end{itemize}

La sintassi per definire un nuovo comando è
\begin{center}
    \texttt{\textbackslash{}newcommand\{nome\_comando\}[argomenti]
    \{definizione\}}
\end{center}

\begin{exampleblock}{Esempio nuovo comando}
	\begin{lstlisting}
\newcommand{\trecolori}[3]{
    \textcolor{red}{#1} 
    \textcolor{green}{#2} 
    \textcolor{blue}{#3}
}
	\end{lstlisting}
\end{exampleblock}


\end{frame}

    \begin{frame}[fragile]{Esercizi comandi personalizzati}

\begin{block}{Esercizio 1 - Concatenazione di stringhe}
    
    \begin{itemize}
        \item [1a] Definire un comando \texttt{\textbackslash myConcat1} che accetti due argomenti e li concateni.
        \item [1b] Definire un comando \texttt{\textbackslash myConcat2} che accetti due argomenti di cui il primo opzionale e li concateni (se il primo non è dato mostra solo il secondo).
    \end{itemize}

\end{block}

\begin{block}{Esempi}

    \begin{itemize}
        \item \texttt{\textbackslash myConcat1}\{Yin\}\{Yang\} = YinYang
        \item \texttt{\textbackslash myConcat2}\{Yang\} = Yang
    \end{itemize}

\end{block}

\begin{figure}[l]
    \includegraphics[scale=0.15, left]{res/images/yinyang.jpg}
\end{figure}


\end{frame}

    \begin{frame}[fragile]{Esercizi comandi personalizzati}

\begin{block}{Esercizio 2 - Wikipedia command}
    
Definire un nuovo comando \texttt{\textbackslash wiki} che completi un url di wikipedia.

Il comando deve accettare due parametri; il primo è il testo da visualizzare al posto dell'url ed è opzionale, il secondo è la sezione di testo finale dell'url, cioè quella che ne contraddistingue il contenuto.

\end{block}

\begin{block}{Esempi}

    \begin{itemize}
        \item \texttt{\textbackslash myWiki\{Latex\}} = \myWiki{Latex}
        \item \texttt{\textbackslash myWiki[l'autore di Latex]\{Leslie\_Lamport\}} = \myWiki[l'autore di Latex]{Leslie_Lamport}
    \end{itemize}

\end{block}

\begin{figure}[l]
    \includegraphics[scale=0.03, left]{res/images/wiki.png}
\end{figure}

\end{frame}

    \begin{frame}[fragile]{Soluzioni degli esercizi}

\begin{block}{Esercizio 1a - Concatenazione di stringhe}

\begin{lstlisting}
\newcommand{\myConcatOne}[2] {
	#1#2
}
\end{lstlisting}

\end{block}

\begin{block}{Esercizio 1b - Concatenazione di stringhe}

\begin{lstlisting}
\newcommand{\myConcatTwo}[2][] {
	{\ifthenelse{\equal{#1}{}}{#2}{#1#2}}
}
\end{lstlisting}

\end{block}

\end{frame}

    \begin{frame}[fragile]{Soluzioni degli esercizi}

\begin{block}{Esercizio 2 - Wikipedia command}

\begin{lstlisting}
\newcommand{\myWiki}[2][] {
	\href{http://en.wikipedia.org/wiki/#2}
	{\ifthenelse{\equal{#1}{}}{http://en.wikipedia.org/wiki/#2}{#1}}
}
\end{lstlisting}

\end{block}

\end{frame}


\end{document}
